\documentclass{article}
\input{../../../../lib.sty}

\title{\texttt{fn match\_datetime\_component}}
\author{Michael Shoemate}
\date{}

\begin{document}

\maketitle

\contrib
Proves soundness of \rustdoc{transformations/make\_stable\_expr/namespace\_dt/expr\_datetime\_component/fn}{match\_datetime\_component} in \asOfCommit{mod.rs}{f5bb719}.

\texttt{match\_datetime\_component} returns the data type and number of possible unique values of a compute operation that retrieves a datetime component.

\section{Hoare Triple}
\subsection*{Precondition}
\subsubsection*{Compiler-verified}
\begin{itemize}
    \item Argument \texttt{temporal\_function} of type \texttt{TemporalFunction} 
\end{itemize}

\subsubsection*{Caller-verified}
None

\subsection*{Pseudocode}
\lstinputlisting[language=Python,firstline=2,escapechar=`]{./pseudocode/match_datetime_component.py}

\subsection*{Postcondition}
\validTransformation{\texttt{(temporal\_function)}}{\texttt{match\_datetime\_component}}

\section{Proof}

\begin{proof}
    For each compute operation, the data type is automatically tested and verified to be correct.
    Only output cardinalities that are specified must be proven.
    Of note:
    \begin{itemize}
        \item there may be up to 53 weeks in a year, due to partial weeks overlapping the end of the year.
        \item There are up to 366 ordinal days due to leap years.
        \item Polars does not account for leap seconds, so there are at most 60 seconds in a minute.
        \item Milliseconds, microseconds and nanoseconds are computed modulo seconds.
    \end{itemize}


\end{proof}

\end{document}
